\section*{Chapter \ref{ch:b3:wage-structure} --- The Wage Structure}

\begin{solution}{exo:b3:wage-structure:1}
1979--1999, at 1.19 log points a year. Both sides contributed: relative demand grew faster
than in any other period (4.21~\% a year) and supply growth slowed from 3.51 to 2.28~\%. % figdata: race
\end{solution}

\begin{solution}{exo:b3:wage-structure:2}
$1{,}200/(2/10{,}000) = 6$ million dalens.
\end{solution}

\begin{solution}{exo:b3:wage-structure:3}
Explained $0.08 \times (14 - 12) = 0.16$; unexplained $0.25 - 0.16 = 0.09$.
\end{solution}

\begin{solution}{exo:b3:wage-structure:4}
Gaps that do not shrink when the résumé carries more information about productivity
--- experience, references, skills --- since a statistical discriminator should rely
less on the name as information grows.
\end{solution}

\begin{solution}{exo:b3:wage-structure:5}
$-0.05 = (0.20 - 0.30)/\sigma$, so $\sigma = 2$.
\end{solution}

\begin{solution}{exo:b3:wage-structure:6}
The tested solution is \code{ex06_oaxaca.py}, in the chapter's \code{exercises/}
folder. With white coefficients the explained 0.086 splits into schooling 0.048, age 0.067,
age squared $-0.056$ and marriage 0.028; with black coefficients, 0.107 into 0.062, 0.098, $-0.081$
and 0.029. % output: ex06_oaxaca
Schooling explains most: the age terms largely offset each other.
\end{solution}

\begin{solution}{exo:b3:wage-structure:7}
The premium falls by $0.10/1.5 = 0.067$ log points, about 6.7~\%.
\end{solution}

\begin{solution}{exo:b3:wage-structure:8}
Surplus $\tfrac12 \times 0.6 \times 0.3 \times 0.1^2 = 0.0009$: 0.09~\% of national income. The wage falls by about $0.3 \times 0.1 = 3~\%$, so
native workers, whose wage bill is about $0.6 \times 0.9 = 0.54$ of income, transfer about $0.03 \times 0.54 = 1.6~\%$ of
national income to the owners of other factors: the transfer is nearly twenty times the
surplus.
\end{solution}

\begin{solution}{exo:b3:wage-structure:9}
Workers who move into union jobs earning, before they move, as much as those who
remain non-union earn after; or no change in a worker's wage when he enters or leaves
a union job, once errors in reported status are corrected.
\end{solution}

\begin{solution}{exo:b3:wage-structure:10}
The derivation is in the text. $D_t = \sigma\ln\tfrac{\alpha}{1-\alpha} + (\sigma - 1)\ln(A_s/A_u)$ rises with $A_s/A_u$ if and only if $\sigma > 1$.
\end{solution}

\begin{solution}{exo:b3:wage-structure:11}
The prejudiced employer earns $20 - 20 = 0$ per worker; to him group B costs $16 \times 1.25 = 20$, the same.
The unprejudiced rival hiring group B earns $20 - 16 = 4$ per worker. With free entry such
rivals multiply, bidding up group B's wage until the gap, and their extra profit,
disappear.
\end{solution}

\begin{solution}{exo:b3:wage-structure:12}
The tested solution is \code{ex12_mariel.py}, in the chapter's \code{exercises/} folder.
On raw means the difference in differences is $-0.094$, within kinds $-0.027$, which is noise
around the true zero. % output: ex12_mariel
The raw estimate records the change in who was sampled, not a change in wages.
\end{solution}

\begin{solution}{pb:b3:wage-structure:graduates}
% figdata: brenland_series, brenland_stats
\textbf{1.} 0.382, 0.234 and 0.350.

\textbf{2.} $-1.196$, $-0.513$ and $-0.120$.

\textbf{3.} $(-0.513 + 1.196)/15 = 4.56~\%$ a year to 2005; $(-0.120 + 0.513)/20 = 1.96~\%$ after.

\textbf{4.} The premium fell to 2005 and rose after. Supply grew in both periods; a rising
premium with rising supply needs demand growing faster than supply.

\textbf{5.} That demand grows as fast as supply.

\textbf{6.} $-0.562$ (s.e.\ 0.021).

\textbf{7.} $\sigma = 1/0.562 = 1.78$.

\textbf{8.} The trend coefficient times $\sigma$: $0.0165 \times 1.78 = 2.94~\%$ a year.

\textbf{9.} Above the 1.62 estimated for the United States, inside the range of 1 to 2.5.

\textbf{10.} Demand grows over time and so does supply: without the trend, the supply
coefficient would pick up the demand shift.

\textbf{11.} Errors in measuring supply, or supply responding to the premium itself.

\textbf{12.} $-0.120 + 0.02 \times 10 = 0.080$.

\textbf{13.} The fitted premium in 2035 is 0.399.

\textbf{14.} $e^{0.399} = 1.49$: graduates earn 49~\% more.

\textbf{15.} At the rate demand grows, 2.94~\% a year.

\textbf{16.} Demand growth may slow, as it did in the United States after 1992.

\textbf{17.} $0.08 \times (13 - 12) + 0.02 \times (18 - 20) = 0.04$.

\textbf{18.} $0.15 - 0.04 = 0.11$.

\textbf{19.} Experience raises wages, and B has more of it: on that count B should earn more.

\textbf{20.} Group B's own regression of wages on the same characteristics.

\textbf{21.} Hours, occupation, field of study, firm size.

\textbf{22.} It contains whatever the characteristics miss, and the characteristics may
themselves carry the effects of discrimination.

\textbf{23.} A \omterm{def:b3:wage-structure:audit}{correspondence study}: identical applications differing only in the group
signal.

\textbf{24.} Add a graduate indicator to both regressions: if A has more graduates,
part of the unexplained gap moves into the explained part.

\textbf{25.} $\sigma = 1.78$; the forecast premium in 2035 is 0.40 log points, a wage ratio of 1.49.
\end{solution}
